\newpage
\phantomsection
\label{prf:scalar-mult-positive}
\LRAProofFor{thm:scalar-mult-positive}

\begin{remark*}[Return]
\hyperref[thm:scalar-mult-positive]{Return to Theorem}
\end{remark*}

\ProofVaultURL{https://github.com/wsollers/lra-proof-vault/tree/master/volume-iii/book-analysis-i/bounding/thm-scalar-mult-positive}

\begin{theorem*}[Scalar Multiplication by a Positive Scalar]
Let $A\subseteq\mathbb{R}$ be nonempty and bounded above, and let $k>0$. Then
\[
  \sup(kA)=k\sup A,
  \qquad
  kA=\{ka:a\in A\}.
\]
\end{theorem*}

\begin{proof}[Professional Standard Proof]
\LRAProofBodyStart
Let $s=\sup A$. Since $s$ is an upper bound of $A$, we have $a\le s$ for
every $a\in A$. Because $k>0$, multiplication by $k$ preserves order, so
$ka\le ks$ for every $a\in A$. Thus $ks$ is an upper bound of $kA$.

Let $t=\sup(kA)$. Since $ks$ is an upper bound of $kA$, the least-upper-bound
property gives $t\le ks$. Conversely, since $t$ is an upper bound of $kA$,
we have $ka\le t$ for every $a\in A$. Dividing by $k>0$ gives
$a\le t/k$ for every $a\in A$, so $t/k$ is an upper bound of $A$. Since
$s=\sup A$, we get $s\le t/k$, and multiplying by $k>0$ gives $ks\le t$.
Hence $t=ks$, so $\sup(kA)=k\sup A$.
\end{proof}

\begin{proof}[Detailed Learning Proof]
\LRAProofBodyStart
Set $s=\sup A$. We prove that $ks$ is exactly the least upper bound of
$kA$.

First we show that $ks$ is an upper bound of $kA$. Let $y\in kA$. By the
definition of $kA$, there is some $a\in A$ such that $y=ka$. Since
$s=\sup A$, every element of $A$ is at most $s$, so $a\le s$. Because
$k>0$, multiplication by $k$ preserves the inequality, and therefore
$ka\le ks$. Hence $y\le ks$. Since $y\in kA$ was arbitrary, $ks$ is an
upper bound of $kA$.

Now let $t=\sup(kA)$. From the preceding paragraph, $ks$ is an upper bound
of $kA$, so the defining leastness of $t$ gives
\[
  t\le ks.
\]
It remains to prove the reverse inequality. Since $t$ is an upper bound of
$kA$, for each $a\in A$ the element $ka$ belongs to $kA$ and therefore
$ka\le t$. Because $k>0$, dividing by $k$ preserves order, so
\[
  a\le \frac{t}{k}
\]
for every $a\in A$. Thus $t/k$ is an upper bound of $A$. Since $s=\sup A$,
$s$ is less than or equal to every upper bound of $A$, so
\[
  s\le \frac{t}{k}.
\]
Multiplying by $k>0$ gives $ks\le t$.

Together, $t\le ks$ and $ks\le t$ imply $t=ks$. Since
$t=\sup(kA)$ and $s=\sup A$, this is precisely
\[
  \sup(kA)=k\sup A.
\]
\end{proof}

\begin{remark*}[Proof structure]
Show that the proposed value $k\sup A$ is an upper bound of $kA$, then use
the least-upper-bound property in the reverse direction by dividing an
arbitrary upper bound of $kA$ by the positive scalar $k$.
\end{remark*}

\begin{dependencies}
\begin{itemize}
  \item \hyperref[def:real-upper-bound]{Upper bound}.
  \item \hyperref[def:supremum]{Supremum}.
  \item \hyperref[prop:upper-bound-comparison-with-the-supremum]{Upper Bound Comparison with the Supremum}.
  \item \hyperref[thm:ineq-nonstrict-multiply-positive]{Multiplication by a positive scalar preserves non-strict order}.
\end{itemize}
\end{dependencies}

\clearpage
